% from ega4-1-fr.tex
\oldpage[IV-1]{231}

  \item[(iii)] Si $f:X\to Y$ est un $S$-morphisme localement de présentation
  finie, $f_{(S')}:X_{(S')}\to Y_{(S')}$ est localement de présentation finie
  pour toute extension $S'\to S$ du préschéma de base.
  \item[(iv)] Si $f:X\to X'$ et $g:Y\to Y'$ sont deux $S$-morphismes
  localement de présentation finie, $f\times_Sg$ est localement de
  présentation finie.
  \item[(v)] Si le composé $g\circ f$ de deux morphismes est localement de
  présentation finie et si $g$ est localement de type fini, alors $f$ est
  localement de présentation finie.
\end{enumerate}
\end{proposition}

\begin{proof}
L'assertion (i) est triviale. Pour démontrer (iii), on peut se limiter au cas
où $S=Y$ (I, 3.3.11)~; soient $x'$ un point de $X_{(S')}$, $s'$ et $x$ ses
projections sur $S'$ et $X$ respectivement, $s$ la projection commune de $s'$
et $x$ sur $S$. Par hypothèse, il y a un voisinage ouvert affine $V$ (resp.
$U$) de $s$ (resp. $x$) tels que l'image de $U$ soit contenue dans $V$ et que
$A(U)$ soit une $A(V)$-algèbre de présentation finie~; soit $V'$ un voisinage
ouvert affine de $s'$ dont l'image dans $S$ soit contenue dans $V$~; alors
$U\times_VV'$ est un voisinage ouvert affine de $x'$ dont l'image dans $S'$
est contenue dans $V'$ et dont l'anneau $A(U)\otimes_{A(V)}A(V')$ est une
$A(V')$-algèbre de présentation finie (1.4.1).

Pour démontrer (ii), considérons un point $x\in X$, et soient $y=f(x)$,
$z=g(f(x))$. Il y a par hypothèse un voisinage ouvert affine $W=\operatorname{Spec}(C)$
(resp. $V=\operatorname{Spec}(B)$) de $z$ (resp. $y$) tels que $g(V)\subset W$ et que $B$
soit une $C$-algèbre de présentation finie. D'autre part, il y a un voisinage
ouvert affine $V'=\operatorname{Spec}(B')$ (resp. $U=\operatorname{Spec}(A)$) de $y$ (resp. $x$) tels que
$f(U)\subset V'$ et que $A$ soit une $B'$-algèbre de présentation finie. Soit
$s\in B$ tel que l'ouvert $D(s)=\operatorname{Spec}(B_s)$ dans $V$ soit un voisinage de $y$
contenu dans $V'$, et soit $U'=f^{-1}(D(s))\subset U$~; on a
$U'=\operatorname{Spec}(A\otimes_{B'}B_s)$, et $A\otimes_{B'}B_s$ est une $B_s$-algèbre de
présentation finie (1.4.1)~; d'aute part $B_s=B[T]/(1-sT)$ est une
$B$-algèbre de présentation finie par définition~; par suite (1.4.1),
$A\otimes_{B'}B_s$ est une $C$-algèbre de présentation finie, ce qui prouve
(ii).

On sait que (iv) résulte de (i), (ii) et (iii) (I, 3.5.1). Pour prouver (v),
considérons le diagramme commutatif
\[
\xymatrix{
X \ar[r]^{\Gamma_f} \ar[d] & X\times_ZY \ar[d]^{f\times 1} \\
Y \ar[r]_{\Delta_g} & Y\times_ZY
}
\]
qui identifie $X$ au produit des $(Y\times_ZY)$-préschémas $Y$ et
$X\times_ZY$ (I, 5.3.7, $\Delta_g$ étant le morphisme diagonal). Si l'on sait
que $\Delta_g$ est localement de présentation finie, il en résultera par
(iii) qu'il en est de même de $\Gamma_f$. D'autre part, on a la factorisation
de $f$ en
$X\xrightarrow{\Gamma_f}X\times_ZY\xrightarrow{p_2}Y$ (I, 5.3.13), où la
projection $p_2$ est égale à $(g\circ f)\times_Z1_Y$~; comme $g\circ f$ est
supposé localement de présentation finie, il en est de même de $p_2$ par
(iv), et il résultera finalement de (ii) que $f$ est aussi localement de
présentation finie. On voit donc qu'on est ramené à démontrer le
\end{proof}

\begin{corollary}[1.4.3.1]
\label{IV.1.4.3.1-fr}
Soit $g:Y\to Z$ un morphisme localement de type fini~; alors le morphisme
diagonal $\Delta_g:Y\to Y\times_ZY$ est localement de présentation finie.
\end{corollary}

